\documentclass[11pt, a4paper]{article}

% ── Packages ──────────────────────────────────────────────────────
\usepackage[margin=1in]{geometry}
\usepackage{amsmath, amssymb, amsthm}
\usepackage{graphicx}
\usepackage{booktabs}
\usepackage{hyperref}
\usepackage{cleveref}
\usepackage[backend=biber, style=numeric, sorting=nyt]{biblatex}
\usepackage{caption}
\usepackage{subcaption}
\usepackage{float}
\usepackage{microtype}
\usepackage{bm}

\addbibresource{references.bib}
\graphicspath{{../figures/}}

% ── Macros ────────────────────────────────────────────────────────
\newcommand{\R}{\mathbb{R}}
\newcommand{\pd}[2]{\frac{\partial #1}{\partial #2}}
\newcommand{\norm}[1]{\left\lVert #1 \right\rVert}
\newcommand{\vect}[1]{\bm{#1}}

% ── Title ─────────────────────────────────────────────────────────
\title{Self-Supervised Dense Displacement Estimation\\for Particle Image Velocimetry\\Using Iterative Refinement}
\author{Wael Mansour\\[4pt]
\textit{Mathematical Project}\\
\textit{Bachelor's Programme}}
\date{\today}

% ══════════════════════════════════════════════════════════════════
\begin{document}
\maketitle

% ── Abstract ──────────────────────────────────────────────────────
\begin{abstract}
Particle Image Velocimetry (PIV) is a widely used experimental technique for measuring fluid velocity fields by tracking tracer particles between consecutive image frames.
Traditional correlation-based PIV methods discretise the domain into coarse interrogation windows, limiting spatial resolution.
In this project, I develop a deep learning approach that estimates \emph{dense, per-pixel} displacement fields from synthetic particle image pairs without requiring ground-truth flow labels during training.
The method uses a self-supervised loss composed of a photometric reconstruction term, a divergence-free physics constraint, and a spatial smoothness regulariser.
I implement and compare two architectures: a U-Net baseline and a RAFT-style iterative flow estimator based on correlation volumes and a convolutional gated recurrent unit.
The RAFT model achieves a mean end-point error of 0.31\,px on synthetic test flows---a $2.8\times$ improvement over the U-Net baseline (0.87\,px)---demonstrating that iterative refinement with explicit feature matching significantly outperforms single-pass regression for this task.
\end{abstract}

% ══════════════════════════════════════════════════════════════════
\section{Introduction}
\label{sec:intro}

Measuring fluid velocity fields is fundamental to experimental fluid mechanics, climate modelling, biomedical flow analysis, and industrial process control.
Particle Image Velocimetry (PIV) achieves this by seeding a fluid with small tracer particles, illuminating them with a laser sheet, and capturing two consecutive images separated by a known time interval~\cite{adrian2011particle}.
The displacement of each particle between frames, divided by the time step, yields a local velocity measurement.

Classical PIV algorithms divide the image into interrogation windows (typically $32 \times 32$ or $64 \times 64$ pixels) and compute cross-correlations to find the most likely displacement within each window.
While robust, this approach inherently limits spatial resolution to the window size and struggles with strong velocity gradients within a single window.

Optical flow methods from computer vision offer a per-pixel alternative.
The seminal work of Horn and Schunck~\cite{horn1981determining} formulated optical flow estimation as a variational problem: minimise the brightness constancy residual subject to a global smoothness constraint.
Lucas and Kanade~\cite{lucas1981iterative} proposed a local, window-based variant that is more robust to noise but, like classical PIV, sacrifices resolution.
More recently, deep learning approaches---beginning with FlowNet~\cite{dosovitskiy2015flownet} and refined through PWC-Net~\cite{sun2018pwcnet} and RAFT~\cite{teed2020raft}---have achieved state-of-the-art accuracy on natural image benchmarks by learning to predict optical flow end-to-end.

A key challenge in applying deep learning to PIV is the scarcity of labelled data: obtaining pixel-accurate ground-truth displacement fields for real experimental images is impractical.
Self-supervised (also called unsupervised) optical flow methods address this by training with proxy losses that do not require ground truth~\cite{jason2016back, meister2018unflow}.
The core idea is that if the predicted flow is correct, then warping frame~1 by the predicted displacement should reconstruct frame~2.

In this project, I combine self-supervised training with physics-informed regularisation to estimate dense displacement fields for PIV.
I use synthetically generated particle images with parametric flow fields, enabling controlled evaluation against known ground truth.
I implement and compare two neural network architectures:
\begin{enumerate}
    \item A \textbf{U-Net} encoder--decoder baseline~\cite{ronneberger2015unet}, which predicts flow in a single forward pass.
    \item A \textbf{RAFT-style iterative estimator}~\cite{teed2020raft}, which builds a correlation volume between frame features and refines the flow over 12 iterations using a convolutional GRU.
\end{enumerate}
The RAFT model achieves a mean end-point error of 0.31\,px, compared to 0.87\,px for the U-Net---demonstrating the value of iterative refinement and explicit feature matching for particle tracking.

% ══════════════════════════════════════════════════════════════════
\section{Background}
\label{sec:background}

\subsection{Optical Flow and the Brightness Constancy Assumption}

Optical flow is the apparent motion of intensity patterns between two images.
Given a grayscale image $I(\vect{x}, t)$ where $\vect{x} = (x, y)$ denotes spatial position and $t$ denotes time, the \emph{brightness constancy assumption} states that the intensity of a moving point does not change:
\begin{equation}
    I(\vect{x}, t) = I(\vect{x} + \vect{d}, t + \Delta t),
    \label{eq:bca}
\end{equation}
where $\vect{d} = (u, v)$ is the displacement vector.
A first-order Taylor expansion of the right-hand side yields the \emph{optical flow constraint equation}:
\begin{equation}
    \pd{I}{x}\, u + \pd{I}{y}\, v + \pd{I}{t} = 0.
    \label{eq:ofce}
\end{equation}
This is a single equation in two unknowns, making the problem ill-posed without additional constraints.

Horn and Schunck~\cite{horn1981determining} resolved this by introducing a global smoothness regulariser, formulating the estimation as a variational problem:
\begin{equation}
    \min_{\vect{d}} \int_\Omega \left[ \left( \pd{I}{x} u + \pd{I}{y} v + \pd{I}{t} \right)^2 + \alpha \left( \norm{\nabla u}^2 + \norm{\nabla v}^2 \right) \right] d\vect{x},
    \label{eq:horn_schunck}
\end{equation}
where $\alpha > 0$ controls the trade-off between data fidelity and smoothness.
This variational perspective---balancing a data term with a regulariser---directly informs the loss function design in this project.

\subsection{Deep Learning for Optical Flow}

FlowNet~\cite{dosovitskiy2015flownet} was the first end-to-end convolutional neural network for optical flow, introducing the idea of learning flow prediction directly from image pairs.
PWC-Net~\cite{sun2018pwcnet} improved upon this with a coarse-to-fine pyramid, learnable feature warping, and cost volumes.
RAFT~\cite{teed2020raft} represented a paradigm shift: rather than coarse-to-fine refinement across pyramid levels, RAFT operates at a single resolution ($1/8$) and uses \emph{iterative updates} via a recurrent unit to progressively refine the flow estimate.

The key components of RAFT that I adapt for particle tracking are:
\begin{itemize}
    \item \textbf{Correlation volumes.} An all-pairs dot product between feature maps of the two frames yields a 4D tensor encoding the visual similarity between every pair of spatial locations. A multi-scale pyramid built by average-pooling allows capturing both fine and coarse displacements.
    \item \textbf{Iterative refinement.} A convolutional Gated Recurrent Unit (GRU)~\cite{cho2014gru} maintains a hidden state and updates the flow estimate by a learned residual at each iteration. Crucially, the flow gradient is \emph{detached} between iterations, so only the current residual is trained---preventing vanishing gradients through long iteration chains.
    \item \textbf{Convex upsampling.} Rather than bilinear interpolation, RAFT predicts a convex combination of neighbouring low-resolution flow values to upsample from $1/8$ to full resolution, producing sharper flow boundaries.
\end{itemize}

\subsection{Self-Supervised Optical Flow}

Supervised optical flow methods require ground-truth displacement fields, which are expensive or impossible to obtain for real-world images.
Self-supervised approaches replace the supervised loss with a photometric reconstruction objective: warp frame~1 by the predicted flow and penalise the difference with frame~2~\cite{jason2016back, meister2018unflow}.
This approach works because a correct flow field, by definition, maps pixels in frame~1 to their corresponding locations in frame~2.

For PIV applications, I augment the photometric loss with physics-informed constraints. Since many flows of interest (e.g., low-speed liquid flows) satisfy the incompressibility condition $\nabla \cdot \vect{u} = 0$, I include a divergence penalty. This is a concrete example of embedding domain knowledge into a learning framework---a theme central to modern scientific computing.

% ══════════════════════════════════════════════════════════════════
\section{Methodology}
\label{sec:method}

\subsection{Synthetic Data Generation}
\label{sec:data}

All training and evaluation data are generated synthetically, enabling controlled experiments with exact ground truth.
Each sample consists of a pair of $256 \times 256$ grayscale particle images and a corresponding dense displacement field $\vect{d}(x,y) = (u, v) \in \R^{2 \times 256 \times 256}$.

\paragraph{Particle rendering.}
For each sample, $N \in [500, 2000]$ particles are placed at random positions $\{\vect{p}_i\}_{i=1}^{N}$ with random diameters $d_i \in [2, 4]$~pixels.
Each particle is modelled as a 2D Gaussian blob.
Rather than rendering particles one-by-one in a loop, I exploit the separability of the Gaussian kernel: the contribution of particle $i$ to pixel $(x, y)$ is
\begin{equation}
    G_i(x, y) = \exp\!\left(-\frac{(x - p_{i,x})^2}{2\sigma_i^2}\right) \cdot \exp\!\left(-\frac{(y - p_{i,y})^2}{2\sigma_i^2}\right),
    \label{eq:gaussian_sep}
\end{equation}
where $\sigma_i = d_i / 4$.
Defining the matrices $(\vect{G}_x)_{j,i} = \exp\!\big(-\!(j - p_{i,x})^2 / (2\sigma_i^2)\big)$ and $(\vect{G}_y)_{j,i}$ analogously, the full image is
\begin{equation}
    \vect{I} = \vect{G}_y \, \vect{G}_x^\top \in \R^{256 \times 256},
    \label{eq:matmul_render}
\end{equation}
which sums over all $N$ particles via a single matrix multiplication---orders of magnitude faster than a Python loop.
After rendering, Gaussian sensor noise ($\sigma = 0.05$) is added and pixel values are clamped to $[0, 1]$.

\paragraph{Flow fields.}
I use five parametric flow types with randomised parameters to ensure diversity:
\begin{enumerate}
    \item \textbf{Lamb--Oseen vortex:} a classical viscous vortex model~\cite{lamb1932hydrodynamics} with tangential velocity
    \begin{equation}
        v_\theta(r) = \frac{\Gamma}{2\pi r}\left(1 - \exp\!\left(-\frac{r^2}{r_0^2}\right)\right),
        \label{eq:lamb_oseen}
    \end{equation}
    where $\Gamma$ is the circulation, $r_0$ is the core radius, and $r$ is the distance from the vortex centre.
    \item \textbf{Uniform translation:} constant displacement $(u, v) = m(\cos\theta, \sin\theta)$ at a random angle $\theta$ and magnitude $m$.
    \item \textbf{Linear shear:} $u = a(y - 0.5)$, $v = 0$ (or transposed), modelling a linear velocity profile.
    \item \textbf{Channel (Poiseuille) flow:} a parabolic profile $u = a \cdot 4y(1-y)$, $v = 0$, representing fully developed laminar pipe flow.
    \item \textbf{Multi-vortex:} superposition of 2--5 Lamb--Oseen vortices with independent random parameters and signs.
\end{enumerate}
All flow magnitudes are bounded by a maximum displacement of 8~pixels ($\approx 3\%$ of the image width).
Frame~2 is generated by displacing each particle according to the flow field at its position in frame~1, then re-rendering.

\subsection{Self-Supervised Loss Function}
\label{sec:loss}

The total loss combines three terms, none of which requires ground-truth flow.

\paragraph{Photometric loss.}
Given the predicted flow $\hat{\vect{d}}$, frame~1 is warped via differentiable bilinear sampling (backward warping):
\begin{equation}
    \tilde{I}_1(\vect{q}) = I_1\!\left(\vect{q} - \hat{\vect{d}}(\vect{q})\right),
    \label{eq:warp}
\end{equation}
where $\vect{q}$ is a pixel coordinate in frame~2 and the sampling uses \texttt{grid\_sample} with bilinear interpolation.
Note the subtraction: the flow is defined \emph{forward} (frame~1 $\to$ frame~2), so backward warping requires $\vect{q} - \hat{\vect{d}}$.

The photometric loss blends two complementary terms:
\begin{equation}
    \mathcal{L}_\text{photo} = (1 - \alpha)\, \mathcal{L}_\text{Charb} + \alpha\, \mathcal{L}_\text{SSIM},
    \label{eq:photo_blend}
\end{equation}
with $\alpha = 0.5$.
The Charbonnier penalty~\cite{charbonnier1994two}
\begin{equation}
    \mathcal{L}_\text{Charb} = \frac{1}{|\Omega|} \sum_{\vect{q} \in \Omega} \sqrt{\left(\tilde{I}_1(\vect{q}) - I_2(\vect{q})\right)^2 + \epsilon}
    \label{eq:charbonnier}
\end{equation}
(with $\epsilon = 10^{-6}$) is a robust approximation to $L^1$ that is differentiable at zero.
The Structural Similarity (SSIM) loss~\cite{wang2004image}
\begin{equation}
    \mathcal{L}_\text{SSIM} = 1 - \text{SSIM}(\tilde{I}_1, I_2)
    \label{eq:ssim}
\end{equation}
compares local luminance, contrast, and correlation within a Gaussian-weighted window, making it sensitive to \emph{structural} differences rather than raw pixel values.
This is important for particle images where many particles appear visually similar: pixel-wise losses alone can converge to local minima that match the wrong particle, while SSIM enforces consistency of the local neighbourhood structure.

\paragraph{Divergence loss.}
For incompressible flows, the velocity field must satisfy the continuity equation $\nabla \cdot \vect{u} = 0$, or equivalently:
\begin{equation}
    \pd{u}{x} + \pd{v}{y} = 0.
    \label{eq:divergence}
\end{equation}
I penalise deviations from this constraint using central finite differences on the interior pixels:
\begin{equation}
    \mathcal{L}_\text{div} = \frac{1}{|\Omega'|} \sum_{\vect{q} \in \Omega'} \left( \pd{\hat{u}}{x}\bigg|_{\vect{q}} + \pd{\hat{v}}{y}\bigg|_{\vect{q}} \right)^2,
    \label{eq:div_loss}
\end{equation}
where $\Omega'$ is the set of interior pixels (excluding boundaries where central differences are undefined).
This physics prior prevents the model from predicting non-physical flows containing sources or sinks.

\paragraph{Smoothness loss.}
A total-variation-style regulariser penalises large spatial gradients of the flow:
\begin{equation}
    \mathcal{L}_\text{smooth} = \frac{1}{|\Omega|} \sum_{\vect{q}} \left[ \left(\pd{\hat{\vect{d}}}{x}\right)^2 + \left(\pd{\hat{\vect{d}}}{y}\right)^2 \right]_{\vect{q}}.
    \label{eq:smooth_loss}
\end{equation}
This prevents noisy, spatially incoherent flow predictions that could ``game'' the photometric loss by matching individual particles through irregular displacements.
The connection to the Horn--Schunck regulariser in \cref{eq:horn_schunck} is direct: both penalise $\|\nabla u\|^2 + \|\nabla v\|^2$.

\paragraph{Iterative sequence loss.}
For the RAFT model, which produces a flow prediction $\hat{\vect{d}}^{(i)}$ at each iteration $i = 1, \ldots, N$ (with $N=12$), the losses are summed with exponentially increasing weights so that later (more refined) predictions contribute more:
\begin{equation}
    \mathcal{L} = \sum_{i=1}^{N} \gamma^{N-i} \left[ \lambda_\text{photo}\, \mathcal{L}_\text{photo}^{(i)} + \lambda_\text{div}\, \mathcal{L}_\text{div}^{(i)} + \lambda_\text{smooth}\, \mathcal{L}_\text{smooth}^{(i)} \right],
    \label{eq:total_loss}
\end{equation}
with $\gamma = 0.8$, $\lambda_\text{photo} = 1.0$, $\lambda_\text{div} = 0.3$, and $\lambda_\text{smooth} = 0.15$.
The final iteration has weight~1.0, while the first has weight $0.8^{11} \approx 0.09$.

\subsection{Model Architecture}
\label{sec:architecture}

\paragraph{U-Net baseline.}
The baseline is a 4-level U-Net encoder--decoder~\cite{ronneberger2015unet} with skip connections.
It takes the two-frame input $[B, 2, 256, 256]$ and outputs a single flow prediction $[B, 2, 256, 256]$ in one forward pass.
With 7.8M trainable parameters, the U-Net has no mechanism for explicit feature matching between frames---it must learn to compare them implicitly through convolutions.

\paragraph{RAFT-style iterative estimator.}
The main model, \texttt{FlowEstimatorRAFT}, adapts the RAFT architecture~\cite{teed2020raft} to single-channel particle images.
It has 7.1M parameters and operates as follows:

\begin{enumerate}
    \item \textbf{Feature extraction.}
    A shared \emph{FeatureEncoder} (residual CNN with InstanceNorm) processes each frame independently to produce feature maps at $1/8$ resolution: $\vect{F}_k \in \R^{B \times 256 \times 32 \times 32}$ for $k \in \{1, 2\}$.
    A separate \emph{ContextEncoder} processes frame~1 to produce a GRU hidden state $\vect{h}_0 \in \R^{B \times 128 \times 32 \times 32}$ and context features $\vect{c} \in \R^{B \times 128 \times 32 \times 32}$.

    \item \textbf{Correlation volume.}
    The all-pairs correlation between the two feature maps is computed as:
    \begin{equation}
        C_{ij} = \frac{1}{\sqrt{D}} \sum_{d=1}^{D} F_{1}^{(d)}(i)\, F_{2}^{(d)}(j),
        \label{eq:corr}
    \end{equation}
    where $D = 256$ is the feature dimension and $i, j$ index spatial locations.
    This 4D volume is pooled into a 4-level pyramid (successive $2\times$ average pooling), enabling the model to match features at multiple scales.
    At each iteration, a local $(2r+1)^2$ neighbourhood (with $r=4$, so 81 values) is sampled from each pyramid level at the current flow estimate, yielding $4 \times 81 = 324$ correlation features per pixel.

    \item \textbf{Iterative refinement.}
    A \emph{MotionEncoder} fuses the correlation lookup with the current flow estimate into motion features.
    These, combined with the context features $\vect{c}$, are fed to a \emph{convolutional GRU}~\cite{cho2014gru} that updates the hidden state:
    \begin{equation}
        \vect{h}_{t} = (1 - \vect{z}_t) \odot \vect{h}_{t-1} + \vect{z}_t \odot \tilde{\vect{h}}_t,
        \label{eq:gru}
    \end{equation}
    where $\vect{z}_t$ is the update gate and $\tilde{\vect{h}}_t$ is the candidate state.
    A \emph{FlowHead} (two convolutional layers) predicts a residual $\Delta\vect{d}^{(t)}$ from the hidden state, and the flow is updated: $\hat{\vect{d}}^{(t)} = \hat{\vect{d}}^{(t-1)} + \Delta\vect{d}^{(t)}$.

    \item \textbf{Convex upsampling.}
    A \emph{MaskHead} predicts a $8 \times 8 \times 9$ weight mask from the hidden state.
    The $1/8$-resolution flow is upsampled $8\times$ by taking a softmax-weighted convex combination of each pixel's $3 \times 3$ neighbourhood:
    \begin{equation}
        \hat{\vect{d}}_\text{HR}(p) = \sum_{k=1}^{9} w_k(p)\, \hat{\vect{d}}_\text{LR}(\text{nb}_k(p)),
        \label{eq:convex_up}
    \end{equation}
    where $w_k$ are the softmax weights. This produces sharper flow boundaries than bilinear interpolation.
\end{enumerate}

\subsection{Training Details}

Both models are trained for 80 epochs on 4{,}000 synthetically generated samples per epoch, with 400 validation samples.
The RAFT model uses AdamW~\cite{loshchilov2019adamw} with learning rate $2 \times 10^{-4}$ and weight decay $10^{-5}$, with cosine annealing reducing the learning rate to $1\%$ of its initial value.
Automatic mixed precision (AMP)~\cite{micikevicius2018mixed} is used for memory efficiency, and gradient norms are clipped at 1.0 to stabilise training.

% ══════════════════════════════════════════════════════════════════
\section{Experiments and Results}
\label{sec:results}

\subsection{Evaluation Metrics}

The primary metric is the \textbf{End-Point Error (EPE)}: the mean Euclidean distance between predicted and ground-truth displacement vectors:
\begin{equation}
    \text{EPE} = \frac{1}{|\Omega|} \sum_{\vect{q} \in \Omega} \sqrt{(\hat{u}(\vect{q}) - u^*(\vect{q}))^2 + (\hat{v}(\vect{q}) - v^*(\vect{q}))^2}.
    \label{eq:epe}
\end{equation}
Note that EPE is computed only during evaluation---it is never used as a training signal.
I also report the mean divergence $\|\nabla \cdot \hat{\vect{d}}\|^2$ and smoothness $\|\nabla \hat{\vect{d}}\|^2$ of the predicted flow.

\subsection{U-Net vs.\ RAFT Comparison}

\Cref{tab:comparison} summarises the final performance of both architectures.

\begin{table}[H]
    \centering
    \caption{Quantitative comparison of the U-Net baseline and RAFT model on synthetic particle images ($256 \times 256$, max displacement 8\,px).}
    \label{tab:comparison}
    \begin{tabular}{lccc}
        \toprule
        \textbf{Metric} & \textbf{U-Net} & \textbf{RAFT} & \textbf{Improvement} \\
        \midrule
        End-Point Error (px) & 0.87 & \textbf{0.31} & $2.8\times$ \\
        Divergence           & 0.020 & \textbf{0.002} & $10\times$ \\
        Smoothness           & 0.085 & \textbf{0.006} & $14\times$ \\
        Relative Error       & 10.9\% & \textbf{3.9\%} & --- \\
        Parameters           & 7.8M & 7.1M & --- \\
        \bottomrule
    \end{tabular}
\end{table}

The RAFT model achieves sub-pixel accuracy (0.31\,px EPE) with dramatically lower divergence and smoothness losses, indicating that the predicted flow fields are both more accurate and more physically plausible.
Notably, the RAFT model achieves this with \emph{fewer} parameters than the U-Net, suggesting that the architectural inductive biases (correlation volumes, iterative refinement) are more important than raw model capacity.

\subsection{Qualitative Results}

\Cref{fig:vortex_flow,fig:channel_flow,fig:multi_vortex_flow} show predicted vs.\ ground-truth flow fields for three representative flow types, visualised using HSV colour coding (hue encodes direction, brightness encodes magnitude).

\begin{figure}[H]
    \centering
    \includegraphics[width=0.95\linewidth]{vortex/sample_00_2_flow_hsv.png}
    \caption{Lamb--Oseen vortex flow: predicted (left) vs.\ ground truth (right). The model captures the rotational structure and radial decay of the vortex accurately.}
    \label{fig:vortex_flow}
\end{figure}

\begin{figure}[H]
    \centering
    \includegraphics[width=0.95\linewidth]{channel/sample_00_2_flow_hsv.png}
    \caption{Channel (Poiseuille) flow: predicted (left) vs.\ ground truth (right). The parabolic velocity profile is faithfully recovered, though minor artifacts are visible near the walls.}
    \label{fig:channel_flow}
\end{figure}

\begin{figure}[H]
    \centering
    \includegraphics[width=0.95\linewidth]{multi_vortex/sample_00_2_flow_hsv.png}
    \caption{Multi-vortex flow: predicted (left) vs.\ ground truth (right). The model successfully resolves the superposition of multiple vortices with different circulations and positions.}
    \label{fig:multi_vortex_flow}
\end{figure}

\Cref{fig:quiver} shows a quiver plot comparison for the vortex flow, overlaid on the particle image.
The predicted flow arrows closely match the ground truth in both direction and magnitude.

\begin{figure}[H]
    \centering
    \includegraphics[width=0.95\linewidth]{vortex/sample_00_3_quiver.png}
    \caption{Quiver plot for vortex flow: predicted (left) vs.\ ground truth (right), overlaid on the particle image. Arrow colour indicates magnitude (cyan = low, magenta = high).}
    \label{fig:quiver}
\end{figure}

\Cref{fig:error_div} presents the spatial distribution of end-point error and divergence for a vortex sample.
The error is concentrated near the vortex core where the velocity gradient is steepest, which is consistent with the difficulty of resolving rapid spatial variations at $1/8$ resolution.
The divergence field is nearly zero everywhere, confirming that the divergence loss effectively enforces the incompressibility constraint.

\begin{figure}[H]
    \centering
    \includegraphics[width=0.95\linewidth]{vortex/sample_00_5_error_div.png}
    \caption{End-point error heatmap (left, mean 0.47\,px) and divergence field (right) for a vortex sample. Error is highest at the vortex core; divergence is near zero everywhere.}
    \label{fig:error_div}
\end{figure}

\Cref{fig:frames} shows a typical input particle image pair, illustrating the subtle displacement of particles between frames that the network must detect.

\begin{figure}[H]
    \centering
    \includegraphics[width=0.95\linewidth]{vortex/sample_00_1_frames.png}
    \caption{A synthetic particle image pair (vortex flow). The displacement between frame~1 and frame~2 is at most 8 pixels---difficult to discern by eye.}
    \label{fig:frames}
\end{figure}

\subsection{Development Insights}

A critical bug encountered during development illustrates an important subtlety of image warping.
The flow field is defined in the \emph{forward} direction (frame~1 $\to$ frame~2), but PyTorch's \texttt{grid\_sample} performs \emph{backward} sampling: it asks ``where in the source image should I sample to fill this output pixel?''
The correct warping formula is therefore $\tilde{I}_1(\vect{q}) = I_1(\vect{q} - \hat{\vect{d}}(\vect{q}))$, using subtraction.
An initial implementation using addition caused the model to learn $\hat{\vect{d}} \approx -\vect{d}^*$ to compensate, resulting in a photometric loss that decreased while the EPE diverged.
Correcting this sign error immediately reduced the U-Net EPE from 4.3 to 0.87\,px.

% ══════════════════════════════════════════════════════════════════
\section{Discussion and Conclusion}
\label{sec:conclusion}

This project demonstrates that self-supervised deep learning can achieve sub-pixel accuracy for dense displacement estimation in PIV-like settings, without requiring any ground-truth flow labels.
The key findings are:

\begin{enumerate}
    \item \textbf{Iterative refinement outperforms single-pass regression.}
    The RAFT model's correlation volumes provide an explicit mechanism for feature matching between frames, and the ConvGRU enables progressive correction of the flow estimate.
    This $2.8\times$ improvement over the U-Net comes with fewer parameters, highlighting the importance of architectural design over model capacity.

    \item \textbf{Physics-informed losses improve physical plausibility.}
    The divergence loss reduces the mean divergence by $10\times$ compared to the U-Net, producing flow fields that respect the incompressibility constraint.
    The connection between the smoothness loss and the classical Horn--Schunck regulariser (\cref{eq:horn_schunck}) illustrates how variational principles from applied mathematics can be embedded into neural network training.

    \item \textbf{SSIM is critical for particle images.}
    Blending the Charbonnier loss with SSIM prevents the model from converging to local minima where it matches the wrong particles.
    The structural comparison enforced by SSIM provides a richer gradient signal than pixel-wise losses alone.

    \item \textbf{Warping direction matters.}
    The distinction between forward flow and backward sampling is a mathematical subtlety with major practical consequences, as the sign-error bug demonstrated.
\end{enumerate}

\paragraph{Limitations.}
The current approach trains and evaluates on synthetic data only.
Real PIV images exhibit non-uniform illumination, out-of-plane particle motion, and optical distortions that are not captured by the synthetic generator.
Additionally, the maximum displacement of 8\,px is modest; scaling to larger displacements would require multi-scale correlation volumes or coarse-to-fine strategies.

\paragraph{Future work.}
Possible extensions include:
(i)~domain adaptation techniques to bridge the gap between synthetic and real PIV images;
(ii)~incorporating temporal consistency across more than two frames;
(iii)~extending to 3D volumetric PIV (tomographic PIV) where the displacement field has three components;
and (iv)~testing on established PIV benchmarks to compare with traditional correlation-based methods.

In summary, by combining the iterative refinement paradigm of RAFT with self-supervised, physics-informed losses, this project achieves accurate dense displacement estimation that respects the underlying fluid mechanics---a compelling example of how ideas from applied mathematics and deep learning can be synergistically combined.

% ══════════════════════════════════════════════════════════════════
\printbibliography

\end{document}
